\subsection{Division values and a variable-degree interpolation test}

\begin{proposition}[ARI C.7.12; Proved route test, with an explicit limitation]\label{ari:C.7.12}
Let $E:y^2+y=x^3-38x+90$, $Q=(4,-2)$, and suppose, only in
this statement, that $N(Q)$ is algebraic.  Write
\[
 \delta=19^{1/4},\qquad b=\eta_1/\omega_1,\qquad
 g(z)=\delta\sigma(z)e^{-bz^2/2},\qquad Q=\exp_E(u),
\]
where $u,\omega_1$ are real as in C.7.9.  Thus
$|g(z)|=N(\exp_E z)$ for real $z$.

Every point in the CM division hull of $Q$ has algebraic norm.
There are constants $D,H$, depending on the fixed curve, $Q$ and
the hypothetical algebraic number $N(Q)$, such that, for
$nR=mQ$, $|m|\leq n$, $R\ne O$,
\[
 h(N(R))\leq H,
 \qquad
 [\Q(R,N(R)):\Q]\leq
 \begin{cases}Dn^4,&R\in E(\R),\\2Dn^6,&R\in E(\C).
 \end{cases}
\]
Here $h$ is absolute logarithmic height; the complex line means
the chosen complex embedding of an algebraic point.

Under the same hypothesis, any finite collection of real translates
$g(z+a)$ by division-hull logarithms has function-field
transcendence degree at most three over $\overline\Q$, after
adjoining the elliptic coordinates.  Three translates already
attain this cap.  The analogous holomorphic extensions in one
variable of $N(\exp_E(z+a))^2$ for complex division-hull shifts
have the same cap.  A safe common degree envelope for a fixed
finite family on the real division grids is $O(n^5)$, with
uniformly bounded heights.

These proved degree and height envelopes do not verify the
variable-degree criterion of Gramain--Mignotte--Waldschmidt
\citep[Theorem 2.1]{gramain1986values}.  This is a failure of the
available sufficient certificate, not a proof that the actual
degrees violate every possible application of that theorem.
The arithmetic-germ hypotheses in the criteria of
\citet[Theorem 6.1]{herblot2012meromorphic} and
\citet[Theorem 4.16]{gasbarri2010analytic} also remain unavailable;
the relevant nontorsion germs have a transcendental tangent.
\end{proposition}

\begin{proof}
\emph{Propagation and degrees.}
The full division identity in C.4.2 is
\[
 N(R)^{n^2}=N(nR)\delta^{n^2-1}/|\psi_n(R)|.
\]
For nontorsion $R$, its last factor is a nonzero algebraic number.
Consequently algebraicity propagates through multiplication
and division.  The CM multiplication formula in C.7.6 gives
the same conclusion for the CM division hull.  All nonzero
torsion values were already proved algebraic in C.4.
For example, if $T$ has order dividing $n$ and $P$ is algebraic
and nontorsion, division at $P$ and $P+T$ gives directly
\[
 \left(\frac{N(P+T)}{N(P)}\right)^{n^2}
       =\frac{|\psi_n(P)|}{|\psi_n(P+T)|}\in\overline\Q.
\]
No undetermined harmonic normalization is needed.

Put $K_0=\Q(\delta,N(Q))$ and $D=[K_0:\Q]$.
For every nonzero integer $m$, $N(mQ)\in K_0$, since
$\psi_m(Q)$ is rational and its absolute value is rational.
The fibre of $[n]$ over $mQ$ has $n^2$ points, so
$[K_0(R):K_0]\leq n^2$.
If $R$ is real, $|\psi_n(R)|=\pm\psi_n(R)$; adjoining
$N(R)$ therefore costs at most a further factor $n^2$.
For a complex $R$, adjoin both $R$ and $\bar R$, at cost
at most $n^4$, and use
\[
 N(R)^{2n^2}=
 \frac{N(mQ)^2\delta^{2(n^2-1)}}
      {\psi_n(R)\psi_n(\bar R)}.
\]
This proves the stated degree envelopes when $m\ne0$.
For $m=0$, take the limit in the division identity at a nonzero
$n$-torsion point.  The zero of $\psi_n$ is simple in the
uniformizing coordinate, and $N(z)=\delta|z|+o(|z|)$ at zero;
hence
\[
 N(R)^{n^2}=
       \frac{n\delta^{n^2}}{|\partial_z\psi_n(R)|}.
\]
The derivative is algebraic, since the invariant derivation on
the Weierstrass model is algebraic.  The same degree bounds follow.

\emph{Uniform heights.}
We give the coefficient estimate needed here, so the height
assertion does not rely on an unspecified division-polynomial
bound.  On the fixed integral model, $\psi_n$ is a polynomial
in $x$ for odd $n$, and $(2y+1)$ times a polynomial in $x$
for even $n$.  These polynomials have degree $O(n^2)$.
Choose a fixed circle $|x|=A>1$ avoiding the finite branch values.
On its compact inverse image in $E(\C)$, $N$ has a positive
minimum $m_A$, and $|2y+1|$ has a positive minimum.  Let $M$
be the maximum of $N$ on the whole compact elliptic curve.
The rearranged division identity gives
\[
 |\psi_n(R)|\leq M\delta^{n^2-1}m_A^{-n^2}
       \qquad(|x(R)|=A).
\]
Cauchy's coefficient estimate, and division by $2y+1$ for
even $n$, show that the logarithm of the coefficient
$\ell^1$ norm is $O_E(n^2)$.  The same assertion holds for
$\partial_z\psi_n$, because differentiation and the fixed
Weierstrass equation add only polynomial factors in $n$.

The elementary height estimate for a fixed degree-four
duplication map, followed by telescoping, gives
$h(x(R))=2\widehat h(R)+O_E(1)$ uniformly over
$E(\overline\Q)$.  Since
$\widehat h(R)=(m/n)^2\widehat h(Q)$, the coordinate heights
are uniformly bounded when $|m|\leq n$.
The polynomial evaluation height inequality
\citep[Lemma 1.0]{gramain1986values} now gives
$h(\psi_n(R))=O_{E,Q}(n^2)$, and the same estimate for its
invariant derivative at torsion points.  Also
\[
 h(N(mQ))\leq m^2h(N(Q))+O_{E,Q}(m^2).
\]
Divide the height of the division identity by $n^2$.
For a complex value one uses
$h(|\alpha|)\leq h(\alpha)$, which follows by squaring and
using $|\alpha|^2=\alpha\bar\alpha$.
This proves a uniform bound for $h(N(R))$.

\emph{The translate count is finite in function-field dimension.}
For a real logarithm $a$ of a nonzero division point define
\[
 h_a(z)=\frac{\sigma(z+a)}{\sigma(z)\sigma(a)}e^{-baz}.
\]
Then
\[
 g(z+a)=\delta^{-1}g(z)g(a)h_a(z),\qquad
 h_a(z+\omega_1)=h_a(z),\quad
 h_a(z+\omega_2)=e^{-2\pi ia/\omega_1}h_a(z).
\]
The last two identities are exact sigma automorphy and Legendre's
relation.  A real point in the CM division hull is in the ordinary
division hull: if $nR=\alpha Q$ is real, then
$(\alpha-\bar\alpha)Q=O$, so $\alpha=\bar\alpha\in\Z$.
Indeed a nonzero endomorphism has finite kernel, whereas $Q$
is nontorsion.  The real locus of this negative-discriminant
curve is connected, so every real point has a real logarithm.
Consequently its real division-hull logarithms are precisely
of the form $(mu+k\omega_1)/n$, and lie in
$\Q u+\Q\omega_1$.
For finitely many shifts choose a common denominator
$q$ and write $a_j=(m_ju+k_j\omega_1)/q$; put $v=u/q$.
Then $h_{a_j}/h_v^{m_j}$ is elliptic for
$\Lambda_q=\Z\omega_1+\Z q\omega_2$.
The curve $\C/\Lambda_q$ is isogenous to $E$ and has an
algebraic model with this differential normalization.
Its displayed elliptic functions have algebraic divisors.
Each therefore equals a rational function over $\overline\Q$
times a complex scalar.  Evaluate at a real division point
avoiding their finitely many zeros and poles.  The formula
for $h_a$ in terms of $g$ and the assumed algebraicity of all
real division values make that scalar algebraic too.
Consequently all these translates are algebraic over
\[
 \overline\Q(E)\bigl(g,h_v\bigr).
\]
This proves the upper cap three over $\overline\Q$, not merely
over $\C$.

The cap is attained.  The multiplier
$e^{-2\pi iv/\omega_1}$ is not a root of unity, since $Q$
is nontorsion.  Applying the shifts $z\mapsto z+m\omega_2$
to a hypothetical polynomial relation shows that $h_v$ is
transcendental over the elliptic function field: at a fixed
generic $z$, a nonzero polynomial cannot vanish at infinitely
many distinct powers of this multiplier.  The same shifts
give $|g(z+m\omega_2)|=\exp(c m^2+O(m))$ with $c>0$.
After clearing denominators, any nonzero highest coefficient in
a polynomial in $g$, with coefficients in the elliptic field and
$h_v$, stays bounded
away from zero on a subsequence of the irrational rotation.
The highest power of $g$ then dominates, a contradiction.
Thus $g,h_v$ and the elliptic coordinate have transcendence
degree three.  Finally sigma addition gives
\[
 \wp(z)=\wp(v)-
 \frac{\delta^2g(z+v)g(z-v)}{g(z)^2g(v)^2},
\]
so the three translates $g(z),g(z+v),g(z-v)$ already attain it.

For a complex division shift $a$, the holomorphic extension
from real $z$ of the squared norm is
\[
 F_a(z)=\frac{N(a)^2}{\delta^2}
                  g(z)^2h_a(z)h_{\bar a}(z).
\]
The multiplier of $h_ah_{\bar a}$ depends only on
$a+\bar a\in\Q u+\Q\omega_1$ for this real CM curve.
The preceding finite-index-lattice argument applies again;
evaluating $F_a$ at a real division point makes the remaining
scalar algebraic.  Thus complex shifts of the norm do not
provide arbitrarily many independent functions either.

For completeness, the useful joint degree envelope for the
three generators is $O(n^5)$.  The map
$z\mapsto(\exp_Ez,h_v(z))$ describes the complex analytic
subgroup $\exp_{G_V}(z,(\zeta(v)-bv)z)$ of the algebraic
extension $G_V$, $V=\exp_E(v)$, in the meromorphic algebraic
section coordinate used in C.7.1.  Its restriction to real
$z$ is a compact circle; no compactness or closedness of the
whole complex analytic subgroup is asserted.
Its value at $z=u$ is algebraic, because
$h_v(u)=\delta g(u+v)/(g(u)g(v))$.
Its real period is $\omega_1$.  Hence for
$z=(mu+k\omega_1)/n$ its value is an $n$-division point of
a fixed-field algebraic multiple of that point.
The multiplication map on $G_V$ has degree $n^3$: its kernel
is an extension of $E[n]$ by $\mathbb G_m[n]$.
Adjoining $g(z)$ costs at most another $n^2$.
The fixed isogeny and the finitely many algebraic constants
affect only the implied constant.  Uniform height bounds for
the real translates follow from the same norm-bound proof with
denominator $nq$ and a fixed larger bound on the ratio of
numerator to denominator; quotient expressions give the same
bound for $h_v$.  For fixed complex shifts, use instead the
proved expression for $F_a$ as a rational elliptic function
on a fixed isogeny, times $g^2h_v^m$, with an algebraic scalar.
The heights of the isogenous elliptic coordinates, $g$ and
$h_v$ are bounded, so the rational-function evaluation height
inequality bounds $h(F_a)$ uniformly.  Removable zeros or poles
of this particular expression lie in a fixed finite set and
can be handled by their finitely many actual values.  Taking
the positive square root to recover the complex-shifted norm
divides its height by two and costs at most a fixed factor two
in degree for each of the fixed finite number of shifts.

\emph{The criterion has been tested, not assumed.}
For $k$ analytic functions, GMW Theorem 2.1 uses a point count
$\mu$, joint degree envelope $d$, height envelopes $\psi_j$
and growth envelopes $\theta_j$; its decisive sufficient
inequality is
\[
 d(N)\prod_{j=1}^k
  \max\{d(N)\psi_j(N),\theta_j(R(N))\}
 \leq c_0\mu(N)^{k-1}\varphi(N)^k,
\]
where
\[
 \varphi(N)=\log
 \frac{R(N)^2+r(N)r(N+1)}
      {R(N)(r(N)+r(N+1))}.
\]
We checked this formula and the other hypotheses in the
original printed pages 101--102.  On a fixed analytic disc
with fixed smaller sampling disc, $\varphi$ and the growth
envelopes are constant.  The height envelopes just proved
are constant as well.  Thus the required comparison reduces
to $d^{k+1}\ll\mu^{k-1}$.

At denominator $n$ there are $O(n^2)$ real points with
$nR=mQ$, $0\leq m\leq n$; through denominator $N$ there are
$O(N^3)$.  For $(\wp,g)$ on a disc without a pole, our proved
degree envelope $d=O(N^4)$ leads to the attempted comparison
$N^{12}\ll N^3$.  For three independent translates, the
safe joint envelope $d=O(N^5)$ leads to
$N^{20}\ll N^6$.  Neither comparison is available from
these estimates.  Adding further translates does not repair
this certificate: beyond three, the dependence conclusion
would already be true and would not contradict the norm
hypothesis.  We do not infer an actual degree lower bound
from an upper bound, and do not rule out a sharper arithmetic
estimate or a differently organized interpolation argument.

The alternative variable-degree theorems of Herblot and
Gasbarri require a formal germ over the residue number field,
with additional arithmetic estimates.  At a real nontorsion
division point with coordinate $z$ and algebraic $g(z)$,
\[
 \frac{dg}{dx}=\frac{g(z)}{\wp'(z)}
                         \bigl(\zeta(z)-bz\bigr).
\]
The slope in parentheses is transcendental by the period-line
argument of C.7.9 applied to that nontorsion algebraic point.
The other two factors are nonzero algebraic numbers.
Thus the germ is not even defined over $\overline\Q$;
the required arithmetic-germ hypothesis has not been gained
by passing to division values.
\end{proof}

We also checked the Mahler-functional-equation route against
\citet[Theorem 1.1 and system (1.2)]{adamczewski2017mahler}.
That theorem requires algebraic Taylor coefficients, a linear
power-map functional system and an algebraic regular argument.
The compact normalization here has the transcendental coefficient
$b$, its elliptic logarithm is transcendental, and its duplication
in the algebraic coordinate is a Latt\`es map.  We have not obtained
a representation satisfying the stated hypotheses; we do not
assert that every alternative representation is impossible.
The entropy identity of
\citet[Theorem 5.1]{einsiedler2001entropy} uses a sequence of
transformations built from the denominators of $x(2^nQ)$,
with scale $4^n$.  It does not identify the height with the
entropy of a single finite-dimensional algebraic automorphism,
and supplies no algebraic value for its exponential.

The period-refined algebraic-group rank theorem has also been
checked in the original sources
\citep[Theorem 4.1 and Corollary 4.3]{waldschmidt1988periodrank}
\citep[main theorem, p.~363]{roy1993periodrank}.
Its coefficient is
$\alpha(G)=2\dim G-2d_0(G)-d_1(G)$, where $d_1$
counts multiplicative \emph{quotients}, not the dimension of
the torus subgroup.  For the nonsplit $G_Q$ of C.7.9,
$d_0=d_1=0$: a nonzero character on its torus would kill
a nonzero integer multiple of the nontorsion extension class.
Thus $\alpha(G_Q)=4$.  The line $\ell_c$ has no nonzero
algebraic vector, has period rank $\kappa=1$, and the norm
hypothesis supplies one independent algebraic-point logarithm
in addition to that period.  The theorem allows total rank
$3=(4-1)\dim\ell_c$; the available rank $2$ therefore gives
no contradiction.  The period rank is exactly one because
a second period on this line would, by Legendre's relation,
force $u/\omega_1$ to be rational, contrary to the
nontorsion of $Q$.
